% Unit 113 terjemahan Bahasa Indonesia dari chapter8.tex baris 1229--1352.
% Sumber: Methods of Algebra, Volume 2, commit
% 9a5803ff77dd3257484cb177f851a73770a59dd3 (CC BY 4.0).
% stable-unit-id: o014.aljabr2.chapter8.bar-construction-a
% Cakupan: Bagian 8.7 dari pembukaan hingga konstruksi bar bebas--lupa.

% segment-id: o014.li2.chapter8.bar-construction-a.h001
\section{Konstruksi bar}\label{sec:bar-resolution}
\hypertarget{o014.aljabr2.chapter8.bar-construction-a}{ }
\index{konstruksi bar@konstruksi bar (bar construction)}
% segment-id: o014.li2.chapter8.bar-construction-a.p001
Definisi~\ref{def:monad} menjelaskan apa yang dimaksud dengan monad pada
suatu kategori $\mathcal{C}$, beserta versi dualnya, yaitu komonad. Titik
tolak bagian ini ialah suatu cara untuk menghasilkan objek kosimplisial
teraugmentasi (atau objek simplisial teraugmentasi) dalam
$\End(\mathcal{C})$ dari suatu monad (atau komonad). Dengan cara ini,
banyak konstruksi dasar dalam aljabar homologis dapat dijelaskan secara
alami, termasuk homologi dan kohomologi Hochschild yang akan ditinjau
kembali pada akhir bagian ini (Contoh~\ref{eg:HH-comonad}).

% segment-id: o014.li2.chapter8.bar-construction-a.eg001
\begin{example}\label{eg:monad-simplicial}
	Misalkan $(T,\mu,\eta)$ suatu monad pada kategori $\mathcal{C}$, yakni
	suatu aljabar dalam kategori endofunktor
	$\End(\mathcal{C})=\mathcal{C}^{\mathcal{C}}$. Konstruksi ``aljabar
	berjalan'' dalam Catatan~\ref{rem:walking-algebra} memberikan funktor
	monoidal yang bersesuaian
	% segment-id: o014.li2.chapter8.bar-construction-a.q001
	\[ \cate{FinOrd} \to \End(\mathcal{C}), \quad \mathbf{n} \mapsto T^n , \]
	yaitu objek kosimplisial teraugmentasi dalam $\End(\mathcal{C})$. Secara
	dual, sebuah komonad $(L,\delta,\epsilon)$ pada kategori $\mathcal{D}$
	memberikan objek simplisial teraugmentasi dalam $\End(\mathcal{D})$;
	di sini $\delta:L\to L^2$ dan
	$\epsilon:L\to\identity_{\mathcal{D}}$.

	\begin{itemize}
		\item Dalam bagian ini kita terutama membahas objek simplisial
		teraugmentasi pada $\mathcal{D}$ yang diberikan oleh komonad. Secara
		eksplisit, suku berderajat $n$-nya ialah $L^{n+1}$, sedangkan morfisme
		muka dan degenerasinya mudah dituliskan sebagai
		% segment-id: o014.li2.chapter8.bar-construction-a.q002
		\begin{align*}
			d_i := L^i \epsilon L^{n-i}: L^{n+1} & \to L^n, \\
			s_j := L^j \delta L^{n-j}: L^{n+1} & \to L^{n+2}, \quad 0 \leq i, j \leq n ;
		\end{align*}
		suku berderajat $-1$ dari objek teraugmentasi itu ialah funktor
		identitas. Perhatikan bahwa rumus untuk $d_0$ masih bermakna ketika
		$n=0$ dan memberikan morfisme augmentasi
		$\epsilon:L\to\identity_{\mathcal{D}}$.

		\item Sejalan dengan itu, objek kosimplisial teraugmentasi pada
		$\mathcal{C}$ memiliki $T^{n+1}$ sebagai suku berderajat $n$, dan
		% segment-id: o014.li2.chapter8.bar-construction-a.q003
		\begin{align*}
			d^i & : = T^i \eta T^{n-i}: T^n \to T^{n+1}, \\
			s^j & := T^j \mu T^{n-j}: T^{n+2} \to T^{n+1}, \quad 0 \leq i, j \leq n ;
		\end{align*}
		ketika $n=0$, $d^0$ menjadi morfisme augmentasi
		$\eta:\identity_{\mathcal{C}}\to T$.
	\end{itemize}
\end{example}

% segment-id: o014.li2.chapter8.bar-construction-a.p002
Salah satu sumber penting monad atau komonad ialah pasangan funktor
adjoin. Karena asal-usulnya dalam teori klasik (lihat
\S\ref{sec:HH}), konstruksi-konstruksi semacam ini secara kolektif disebut
\emph{konstruksi bar}.

% segment-id: o014.li2.chapter8.bar-construction-a.eg002
\begin{example}[Konstruksi bar: pasangan adjoin]\label{eg:bar-adjunction}
	Misalkan diberikan pasangan adjoin
	% segment-id: o014.li2.chapter8.bar-construction-a.q004
	\[\begin{tikzcd}
		F : \mathcal{C} \arrow[shift left, r] & \mathcal{D}: G \arrow[shift left, l] ,
	\end{tikzcd} \quad \begin{array}{ll}
		\eta: \identity_{\mathcal{C}} \to GF & \text{unit} \\
		\epsilon: FG \to \identity_{\mathcal{D}} & \text{kounit}
	\end{array}\]
	Menurut Contoh~\ref{eg:adjunction-monad}, dari sini diperoleh
	\begin{itemize}
		\item monad $(T,\mu,\eta):=(GF,G\epsilon F,\eta)$ pada
		$\mathcal{C}$, dengan objek kosimplisial teraugmentasi yang
		bersesuaian dinotasikan oleh
		% segment-id: o014.li2.chapter8.bar-construction-a.q005
		\[ \mathrm{Cobar}(F, G): \cate{FinOrd} \to \End(\mathcal{C}); \]
		\item komonad $(L,\delta,\varepsilon):=(FG,F\eta G,\epsilon)$ pada
		$\mathcal{D}$, dengan objek simplisial teraugmentasi yang
		bersesuaian dinotasikan oleh
		% segment-id: o014.li2.chapter8.bar-construction-a.q006
		\[ \mathrm{Bar}(F, G): \cate{FinOrd}^{\opp} \to \End(\mathcal{D}). \]
	\end{itemize}
	Morfisme muka dan degenerasi yang bersesuaian (beserta versi dualnya)
	telah dijelaskan dalam Contoh~\ref{eg:monad-simplicial}.

	Karena $\mathrm{Bar}(F,G)$ (atau $\mathrm{Cobar}(F,G)$) mengambil nilai
	dalam kategori funktor $\End(\mathcal{D})$ (atau
	$\End(\mathcal{C})$), kita dapat mengomposisikan setiap sukunya dengan
	$G$ di sebelah kiri (atau di sebelah kanan). Komposisi pertama
	menghasilkan objek simplisial teraugmentasi
	$G\mathrm{Bar}(F,G):\cate{FinOrd}^{\opp}\to
	\mathcal{C}^{\mathcal{D}}$, dengan suku umum berupa funktor
	% segment-id: o014.li2.chapter8.bar-construction-a.q007
	\begin{align*}
		G\mathrm{Bar}(F, G)_n & = GL^{n+1} = (GF)^{n+1} G \\
		& = T^{n+1}G: \mathcal{D} \to \mathcal{C}, \quad n \geq -1.
	\end{align*}

	Morfisme muka $d_i:T^{n+1}G\to T^nG$ (dipahami sebagai morfisme
	augmentasi bila $i=n=0$) ialah
	% segment-id: o014.li2.chapter8.bar-construction-a.q008
	\[d_i = \left\{\begin{array}{rll}
		G L^i \epsilon L^{n-i} & = T^i G \epsilon F T^{n-i-1} G & \\
		& = T^i \mu T^{n-i-1} G, & \text{jika}\; 0 \leq i < n, \\
		G L^n \epsilon & = T^n G \epsilon , & \text{jika}\; i = n.
	\end{array}\right.\]

	Morfisme degenerasi dengan sumber suku berderajat $n\geq0$,
	$s_j:T^{n+1}G\to T^{n+2}G$, diberikan oleh
	% segment-id: o014.li2.chapter8.bar-construction-a.q009
	\begin{align*}
		s_j = G L^j \delta L^{n-j} & =
		G L^j F \eta G L^{n-j} \\
		& = T^{j+1} \eta T^{n-j} G, \quad 0 \leq j \leq n.
	\end{align*}

	Objek kosimplisial teraugmentasi $\mathrm{Cobar}(F,G)G$ juga memiliki
	deskripsi yang bersesuaian: suku berderajat $n$-nya tetap
	$T^{n+1}G=GL^{n+1}$, sedangkan
	% segment-id: o014.li2.chapter8.bar-construction-a.q010
	\begin{gather*}
		\left[ d^i: GL^n \to GL^{n+1} \right] =
		\begin{cases}
			GL^{i-1} \delta L^{n-i}, & 0 \leq i < n \\
			\eta GL^n , & i = n,
		\end{cases} \\
		\left[ s^j: GL^{n+2} \to GL^{n+1} \right] =
		GL^j \epsilon L^{n-j+1}, \quad 0 \leq j \leq n.
	\end{gather*}
	Demikian pula, ketika $i=n=0$, $d^0$ harus dipahami sebagai morfisme
	augmentasi $\eta G:G\to GFG=GL$.
\end{example}

% segment-id: o014.li2.chapter8.bar-construction-a.p003
Pembaca pemula disarankan berhenti sejenak di sini. Untuk homomorfisme
aljabar-$\Bbbk$ $f:A\to B$ dan pasangan adjoin
% segment-id: o014.li2.chapter8.bar-construction-a.q011
\[\begin{tikzcd}
	B \dotimes{A} (\cdot): A\dcate{Mod} \arrow[shift left, r] & B\dcate{Mod}: \text{lupa} \arrow[shift left, l]
\end{tikzcd}\]
uraikan nilai $\mathrm{Bar}$ (atau $\mathrm{Cobar}$) pada modul-$B$ kiri
$M$ (atau modul-$A$ kiri $N$), lalu gunakan funktor $\mathrm{C}$ dari
Definisi~\ref{def:unnormalized-chain} untuk menuliskan kompleks rantai
yang bersesuaian. Untuk versi modul kanan, komonad (atau monad) yang
bersesuaian telah dijelaskan sepenuhnya dalam
Contoh~\ref{eg:comonad-change-ring}.

% segment-id: o014.li2.chapter8.bar-construction-a.eg003
\begin{example}[Konstruksi bar: bebas dan lupa]\label{eg:bar-comonad}
	Misalkan $(T,\mu,\eta)$ monad pada kategori $\mathcal{C}$.
	Definisi--Proposisi~\ref{def:free-T-Alg} memberikan pasangan adjoin
	% segment-id: o014.li2.chapter8.bar-construction-a.q012
	\[\begin{tikzcd}
		\mathrm{Free}^T : \mathcal{C} \arrow[shift left, r, "\text{bebas}"] & \mathcal{C}^T: U^T \arrow[shift left, l, "\text{lupa}"] ,
	\end{tikzcd}\]
	dan monad yang diinduksinya persis $(T,\mu,\eta)$. Memasukkan ini ke
	Contoh~\ref{eg:bar-adjunction} menghasilkan objek simplisial
	teraugmentasi
	% segment-id: o014.li2.chapter8.bar-construction-a.q013
	\[ \mathrm{Bar}(T) := \mathrm{Bar}(\mathrm{Free}^T, U^T): \cate{FinOrd}^{\opp} \to \End(\mathcal{C}^T). \]

	Evaluasi pada $(M,a)\in\Obj(\mathcal{C}^T)$ merupakan funktor dari
	$\End(\mathcal{C}^T)$ ke $\mathcal{C}^T$. Dengan mengevaluasi setiap suku
	$\mathrm{Bar}(T)$,
	diperoleh objek simplisial teraugmentasi
	% segment-id: o014.li2.chapter8.bar-construction-a.q014
	\[ \mathrm{Bar}(M, a): \cate{FinOrd}^{\opp} \to \mathcal{C}^T .\]

	Sekarang tinjau $U^T\mathrm{Bar}(M,a):\cate{FinOrd}^{\opp}\to
	\mathcal{C}$; ini juga merupakan hasil mengevaluasi
	$U^T\mathrm{Bar}(T)$ pada $(M,a)$.
	\begin{itemize}
		\item Dengan mensubstitusikan ke Contoh~\ref{eg:bar-adjunction},
		suku umum $U^T\mathrm{Bar}(M,a)$ ialah
		% segment-id: o014.li2.chapter8.bar-construction-a.q015
		\[ U^T \mathrm{Bar}(M, a)_n = T^{n+1} U^T (M, a) = T^{n+1} M, \quad n \geq -1. \]
		\item Tuliskan kembali $\epsilon$ untuk kounit pasangan
		$(\mathrm{Free}^T,U^T)$. Bukti
		Definisi--Proposisi~\ref{def:free-T-Alg} telah menunjukkan bahwa
		% segment-id: o014.li2.chapter8.bar-construction-a.q016
		\[ \left[ \epsilon_{(M, a)} : \mathrm{Free}^T U^T (M, a) = (TM, \mu_M) \to (M, a) \right] = a, \]
		sehingga morfisme muka $U^T\mathrm{Bar}(M,a)$ ialah
		% segment-id: o014.li2.chapter8.bar-construction-a.q017
		\begin{equation*}
			\left[ d_i: T^{n+1}M \to T^n M \right] = \begin{cases}
				T^i \mu_{T^{n-i-1} M}, & 0 \leq i < n \\
				T^n a, & i = n.
			\end{cases}
		\end{equation*}
		Rumus ini juga bermakna ketika $i=n=0$ dan memberikan morfisme
		augmentasi, yakni $a:TM\to M$.
		\item Morfisme degenerasinya ialah
		% segment-id: o014.li2.chapter8.bar-construction-a.q018
		\[ \left[ s_j: T^{n+1} M \to T^{n+2} M \right] =
		T^{j+1} \eta_{T^{n-j} M}, \quad 0 \leq j \leq n .\]
	\end{itemize}

	Dengan demikian, $U^T\mathrm{Bar}(M,a)$ dapat ditulis secara ringkas
	sebagai
	% segment-id: o014.li2.chapter8.bar-construction-a.q019
	\begin{equation}\label{eqn:bar-simplicial-resolution}
		\begin{tikzcd}[column sep=large]
			\cdots \arrow[r, shift left=3, "{\mu_{T^2 M}}"] \arrow[r, shift left=1] \arrow[r, shift right=1] \arrow[r, shift right=3, "{T^3 a}"'] & T^3 M \arrow[r, shift left=2.5, "\mu_{TM}"] \arrow[r, "{T\mu_M}" description] \arrow[r, shift right=2.5, "{T^2 a}"'] & T^2 M \arrow[r, shift left, "{\mu_M}"] \arrow[shift right, r, "Ta"'] & TM \arrow[r, "a"] & M .
		\end{tikzcd}
	\end{equation}

	Di atas telah dijelaskan secara abstrak bahwa anak-anak panah ini memenuhi
	identitas simplisial~\eqref{eqn:simplicial-identity}, walaupun
	verifikasi langsungnya juga tidak mengandung kesulitan mendasar. Selain
	itu, Contoh~\ref{eg:T-split-fork} menunjukkan bahwa ujung kanan diagram
	merealisasikan $a:TM\to M$ sebagai koekualiser.
\end{example}
